\section{Live TLS-record timing and planner-generated cover}\label{supp:live-timing}
This study tests cover on the network records produced by actual dispatch. It compares immediate controller cover, sequential timing matching, concurrent dispatch and a second planner that generates a complete cover execution. The source literature, the earlier receipt-level counterfactuals and this collection retain separate denominators.

\subsection{Population, protocol and execution}
The study reuses the 300 executed-cover fixtures per model, with 120 training, 60 validation and 120 test fixtures, balanced over six requested actions. Qwen3:8B and Llama3.1:8B each execute all five arms, giving 3,000 formal episodes and 120 paired held-out fixtures per cell. Phrasings recur across splits. This is a new timing-layer collection on previously used workloads; earlier cover results were known when the protocol was designed.

LangGraph runs the same business planner against seven local MCP Streamable HTTP services over TLS~1.3/HTTP/2. Model weights and templates are fixed by layer hashes before and after collection. Model options are temperature zero, a 4,096-token context and at most 96 generated tokens per model step. One episode is active at a time; the local inference server has two execution slots so that the shadow chain can run alongside the real chain. Connections are shared by the two chains within an episode. A lock serializes calls to the same service, while calls to different services and planner inference may overlap. The experiment consequently exercises real concurrent chains without claiming unrestricted same-service multiplexing.

Every arm uses a common \path{/mcp} request path and 4,096/16,384-byte request/response body-padding quanta. The real planner receives only its own tool responses. Cover actions execute on isolated state and cannot directly write the real fixture. All requests and terminal responses are journaled; state outcomes are recomputed by replaying the synthetic operations. No episode is replaced after a missing response.

\paragraph{Five arms}
\emph{No cover} supplies the baseline. \emph{Immediate} dispatches the original one or two cover calls directly after their triggering honest response. \emph{Matched, sequential} waits for a sampled gap and holds the server response for a sampled minimum duration, blocking the real chain as required by that schedule. \emph{Matched, concurrent} uses the same draws in a persistent episode-scoped dispatcher while the real planner continues inference. \emph{Planner shadow} runs a complete second episode of the same model on an independently generated ticket, including its own checks, actions and feedback. It shares the service connections but not business state or planner context.

The shadow action is sampled uniformly from all six classes, including the real class. An independent fair draw selects which chain starts first; the later chain begins 0.2--1.5\,s after the first ticket response. The attack is not told which first call is real. Random order removes that particular role cue, while the chain-separation attack tests whether other role information remains.

\paragraph{Historical timing pools}
All no-cover episodes are collected first. Only their 120 training fixtures per model supply honest inter-call gaps and per-destination durations; the pools and collection amendments are frozen and pushed before any formal cover run. Timing-matched arms sample from those pools. A monotonic not-before wait enforces the selected gap, and the server holds a cover response until its selected minimum duration has elapsed. Transport, padding, scheduling and concurrent work still contribute to the achieved record times. The cover-arm order rotates across fixtures. Backend prompt caches persist between arms; a reused real ticket and a new shadow ticket can therefore have different cache histories.

\subsection{Observation and independent timing check}
The primary observer receives destinations, bidirectional ciphertext counts, episode/call windows and relay timestamps at the completion of each encrypted TLS record. It receives no message body, operation, task identifier, ground-truth call role or model-link observation at test inference. Initialization precedes the call windows. TLS~1.3 outer type 23 identifies the encrypted records used to delimit a window; it is not interpreted as an application-plaintext type.

Npcap independently records both ingress directions. TCP sequence reassembly, separated by connection generation, must reproduce every relay record's outer type, length and count. The evaluator then measures relay-versus-packet completion-time error. Relay read completion is a software clock, not an exact NIC arrival timestamp. Supplied call boundaries and destination visibility remain observer privileges; the experiment does not evaluate discovery of those boundaries.

\subsection{Attacks, selection and prespecified decisions}
Size-only candidates are standardized logistic regression, a 300-tree random forest, the frozen no-cover classifier, and last- and penultimate-destination rules. Timing candidates add direct timing-plus-size classification, supervised gap-detector removal, label-free two-component Gaussian-mixture gap removal, context-detector removal and joint decoding of zero, one or two removals. The context detector includes destinations, positions, sizes and within-episode timing relations. LR/RF choices and detector thresholds are selected on validation data; detector selection uses call-role F1 and then precision. Call-role annotations are permitted for training and evaluation, never as test-time inputs.

The shadow-specific attack partitions the observed calls into two causal chains with a beam of 64 candidates and a transition-gap model fitted on no-cover training traffic. A clean-chain classifier predicts each chain's action. A role selector, fitted on labelled shadow training episodes, chooses the real chain using observed chain summaries. It is not given the test roles. Exact partition recovery, real-chain recovery and call-role accuracy are descriptive diagnostics. With perfect recovery of an unordered pair of independent uniform six-class labels, guessing the real label succeeds with probability $7/12$; this is an idealized pair-ambiguity reference, not a general privacy bound.

For each arm and model, validation accuracy selects a size-only attack and an adaptive attack from the complete suite; ties favor the declared size-first order. Test maxima are reported separately as descriptive results and do not replace the frozen choices. The test unit is the paired fixture, with 120 outcomes per cell.

\paragraph{Four testing families}
Family P contains the six adaptive-minus-size comparisons for the sequential, concurrent and shadow arms. Each uses exact paired McNemar testing with Holm correction across the six comparisons. An arm meets the prespecified separability rule if either model gains at least 12/120 with adjusted $p<.05$. Family E tests the same six contrasts with paired t-TOST, margin $\pm12/120$, 119 degrees of freedom and Holm correction. Both models must pass for an arm to meet the equivalence rule. At all ties, the implementation uses the conservative boundary probability $.9^{120}$ instead of zero. Equivalence concerns the validation-selected adaptive attack relative to the selected size attack; it does not establish cover-role indistinguishability or an all-attacker guarantee.

Family R has two immediate-arm replication comparisons: label-free removal against the selected size attack. Replication requires a gain of at least 24/120 and Holm-adjusted $p<.05$ in both models. Family C has eight size-only cover-versus-no-cover comparisons, also exact paired tests with Holm correction. C quantifies the size-view change; remaining adaptive inference is reported alongside it. Ten thousand action-stratified paired fixture bootstrap draws give intervals conditional on the fitted attacks and the reused workload. They omit refitting, attack-selection and workload-selection uncertainty. With 120 pairs, moderate increments or equivalence can remain unresolved.

\subsection{Costs and business outcomes}
Real-chain latency begins at admission before either shadow chain starts and ends when the real planner finishes, including shadow-first waiting but excluding session initialization. Unique whole-connection bytes include setup and teardown without double-counting overlapping windows. Model calls and tokens include both planners in the shadow arm. Requested-goal success, whole-state correctness and honest-path agreement with no cover are reported on the same fixtures. Changes in planner behavior, cache history and shared inference resources accompany these policy comparisons, so an endpoint difference alone cannot isolate one causal mechanism.

\subsection{Complete attack and decision results}
Tables~\ref{tab:live-attacks-qwen}--\ref{tab:live-attacks-llama} give every frozen candidate; Table~\ref{tab:live-tests} gives every test, discordance count and conditional interval. Immediate removal meets R in both models: 112 versus 75 and 108 versus 80, with adjusted $p=1.9\times10^{-8}$ and $4.3\times10^{-6}$. For historical sequential dispatch, P rejects for Llama (104 versus 82; $p=.00078$) but not Qwen (93 versus 76; $p=.0559$). Concurrent dispatch rejects for both (99 versus 66 and 106 versus 68). Each history-based arm consequently meets its prespecified separability rule through at least one model. The context attack uses non-timing inputs as well as time; these comparisons do not isolate the causal contribution of timestamps.

Both shadow selectors choose the same size attack for the size and adaptive roles, so their predictions are identical. Family E therefore passes by construction ($p_{\mathrm{Holm}}=1.94\times10^{-5}$ under the all-tie safeguard). We classify this cell as degenerate: it is reported for completeness and is not counted as evidence that trace distributions or real/cover roles are indistinguishable. Section~\ref{sec:supp-floor} gives the comparison that replaces it. Qwen's last-destination rule wins validation at 35/60, ahead of RF at 33/60, but on test scores 54/120 against RF's 70/120. Llama's selected RF scores 64/120. These selection differences are kept visible instead of choosing a new primary attack after test evaluation. All eight size-view cover-versus-none comparisons in C reject; the adaptive results show why a size-view decrease alone is an incomplete protection assessment.

\begin{table*}[t]\centering\small
\caption{All frozen Qwen3:8B attacks on live-record test observations: correct out of 120 per cell. None is no cover; history-based policies sample minimum gap/response holds. Dashes indicate an attack not applied to no-cover traffic. Validation choices are reported separately from these individual test scores.}
\label{tab:live-attacks-qwen}
\setlength{\tabcolsep}{5pt}
\begin{tabularx}{\textwidth}{@{}Yrrrrr@{}}\toprule
Attack & None & Immediate & History seq. & History concurrent & Shadow\\\midrule
Size/destination LR & 117 & 64 & 61 & 52 & 60\\
Size/destination RF & 116 & 59 & 61 & 52 & 70\\
Frozen no-cover classifier & 117 & 40 & 41 & 40 & 31\\
Last destination & 114 & 75 & 76 & 66 & 54\\
Penultimate destination & 14 & 43 & 42 & 47 & 43\\
Direct timing + size & 104 & 60 & 57 & 47 & 64\\
Supervised gap removal & -- & 112 & 67 & 57 & 20\\
Label-free mixture removal & -- & 112 & 53 & 38 & 34\\
Context removal & -- & 111 & 93 & 87 & 20\\
Joint removal/task decoding & -- & 112 & 99 & 99 & 32\\
Two-chain separation & -- & -- & -- & -- & 49\\
\bottomrule\end{tabularx}
\end{table*}

\begin{table*}[t]\centering\small
\caption{All frozen Llama3.1:8B attacks on live-record test observations: correct out of 120 per cell. None is no cover; history-based policies sample minimum gap/response holds. Dashes indicate an attack not applied to no-cover traffic. Validation choices are reported separately from these individual test scores.}
\label{tab:live-attacks-llama}
\setlength{\tabcolsep}{5pt}
\begin{tabularx}{\textwidth}{@{}Yrrrrr@{}}\toprule
Attack & None & Immediate & History seq. & History concurrent & Shadow\\\midrule
Size/destination LR & 110 & 82 & 79 & 68 & 59\\
Size/destination RF & 112 & 80 & 82 & 68 & 64\\
Frozen no-cover classifier & 110 & 54 & 57 & 57 & 43\\
Last destination & 80 & 65 & 65 & 48 & 43\\
Penultimate destination & 66 & 65 & 64 & 64 & 46\\
Direct timing + size & 110 & 77 & 81 & 75 & 59\\
Supervised gap removal & -- & 108 & 100 & 98 & 20\\
Label-free mixture removal & -- & 108 & 14 & 66 & 43\\
Context removal & -- & 108 & 101 & 103 & 20\\
Joint removal/task decoding & -- & 108 & 104 & 106 & 33\\
Two-chain separation & -- & -- & -- & -- & 34\\
\bottomrule\end{tabularx}
\end{table*}

\begin{table*}[t]\centering\small
\caption{All prespecified live-record test families. P/E compare the validation-selected adaptive and size attacks; R compares label-free immediate removal with size; C compares each cover size attack with no cover. $+$/$-$ count fixture discordances. Intervals are 95\% paired action-stratified bootstrap intervals conditional on the fitted attacks. E uses paired t-TOST with a $\pm10$-point margin; the shadow selectors choose the identical size attack, so all pairs tie. Holm correction is separate within each declared family.}
\label{tab:live-tests}
\setlength{\tabcolsep}{4pt}
\begin{tabularx}{\textwidth}{@{}llYrrrrr@{}}\toprule
Family & Model & Arm & $+$/$-$ & $\Delta$ (pp) & Conditional interval & Raw $p$ & Holm $p$\\\midrule
P & Qwen & History seq. & 32/15 & 14.2 & [3.3, 25.0] & 0.0186 & 0.0559 \\
P & Qwen & History concurrent & 42/9 & 27.5 & [17.5, 38.3] & $3.4\times10^{-6}$ & $1.7\times10^{-5}$ \\
P & Qwen & Shadow & 0/0 & 0.0 & [0.0, 0.0] & 1 & 1 \\
P & Llama & History seq. & 28/6 & 18.3 & [10.0, 26.7] & $2\times10^{-4}$ & $7.8\times10^{-4}$ \\
P & Llama & History concurrent & 41/3 & 31.7 & [22.5, 40.8] & $1.6\times10^{-9}$ & $9.7\times10^{-9}$ \\
P & Llama & Shadow & 0/0 & 0.0 & [0.0, 0.0] & 1 & 1 \\
\midrule
E & Qwen & History seq. & 32/15 & 14.2 & [3.3, 25.0] & 0.7713 & 1 \\
E & Qwen & History concurrent & 42/9 & 27.5 & [17.5, 38.3] & 0.9992 & 1 \\
E & Qwen & Shadow & 0/0 & 0.0 & [0.0, 0.0] & $3.2\times10^{-6}$ & $1.9\times10^{-5}$ \\
E & Llama & History seq. & 28/6 & 18.3 & [10.0, 26.7] & 0.9643 & 1 \\
E & Llama & History concurrent & 41/3 & 31.7 & [22.5, 40.8] & 1 & 1 \\
E & Llama & Shadow & 0/0 & 0.0 & [0.0, 0.0] & $3.2\times10^{-6}$ & $1.9\times10^{-5}$ \\
\midrule
R & Qwen & Immediate & 41/4 & 30.8 & [21.7, 40.0] & $9.3\times10^{-9}$ & $1.9\times10^{-8}$ \\
R & Llama & Immediate & 33/5 & 23.3 & [15.0, 32.5] & $4.3\times10^{-6}$ & $4.3\times10^{-6}$ \\
\midrule
C & Qwen & Immediate & 0/42 & -35.0 & [-43.3, -26.7] & $4.5\times10^{-13}$ & $2.7\times10^{-12}$ \\
C & Qwen & History seq. & 0/41 & -34.2 & [-42.5, -25.8] & $9.1\times10^{-13}$ & $4.5\times10^{-12}$ \\
C & Qwen & History concurrent & 0/51 & -42.5 & [-50.8, -34.2] & $8.9\times10^{-16}$ & $6.2\times10^{-15}$ \\
C & Qwen & Shadow & 0/63 & -52.5 & [-60.8, -43.3] & $2.2\times10^{-19}$ & $1.7\times10^{-18}$ \\
C & Llama & Immediate & 3/33 & -25.0 & [-33.3, -16.7] & $2.3\times10^{-7}$ & $4.5\times10^{-7}$ \\
C & Llama & History seq. & 2/30 & -23.3 & [-31.7, -15.0] & $2.5\times10^{-7}$ & $4.5\times10^{-7}$ \\
C & Llama & History concurrent & 2/44 & -35.0 & [-44.2, -25.8] & $3.1\times10^{-11}$ & $1.2\times10^{-10}$ \\
C & Llama & Shadow & 4/50 & -38.3 & [-47.5, -28.3] & $3.8\times10^{-11}$ & $1.2\times10^{-10}$ \\
\bottomrule\end{tabularx}
\end{table*}


\paragraph{Shadow roles and start order}
The chain attack exactly partitions both chains in 48/120 Qwen and 59/120 Llama fixtures, but exactly selects the real chain in only 29 and 32. Its role assignments are correct on 563/977 and 673/1,255 calls; these calls are dependent within fixtures. Test data contain 53 shadow-first and 67 real-first fixtures per model, with 21 matching real/shadow action labels. The RF scores 28/53 and 42/67 for Qwen, and 25/53 and 39/67 for Llama, respectively. These descriptive strata expose order sensitivity without adding an inferential family. Shadow tasks themselves reach 116/120 and 89/120 requested goals. The fixed beam, role learner, recurring phrasings and different cache histories limit the tested attack suite.

\subsection{Ideal candidate-label reference for planner shadows}\label{sec:supp-floor}
Let the real chain and $k$ shadow chains carry labels $L_0,\dots,L_k$ from $C$ classes, with $L_0$ the real label. An idealized attacker observes the unordered multiset $\{L_0,\dots,L_k\}$, which corresponds to perfect chain separation and per-chain classification, but not the index of the real chain. Under the idealization that the real chain is exchangeable with the shadows given the multiset (start order is randomized and the same planner runs both; reused real fixtures and their cache histories are the known departures), $\Pr[L_0=c\mid\text{multiset}]=n_c/(k+1)$, where $n_c$ counts chains labeled $c$, and the Bayes-optimal guess is a modal label. Its success probability is
\[P^*(k)=\frac{\mathbb{E}[\max_c n_c]}{k+1}.\]
For iid uniform shadow labels and $C=6$, exact enumeration gives $1$, $7/12$, $53/108$, $197/432$ and $185/432$ for $k=0,\dots,4$. If shadow labels are instead drawn without replacement from the $C-1$ classes other than the real one, every multiset has distinct labels, the posterior is uniform over the $k+1$ chains, and $P^*(k)=1/(k+1)$ for $k\le C-1$. This variant keeps each shadow's marginal label uniform. The script \texttt{experiment/shadow\_floor.py} computes these values and writes \texttt{analysis/shadow-floor.json}.

On a realized split with $m$ matching pairs among $n$ fixtures and $k=1$, the ideal reference is $(m+(n-m)/2)/n$. The live test split has $m=21$ and $n=120$, giving 70.5/120. The strongest tested call-window attack scores 70/120 (Qwen) and 64/120 (Llama). This is optimal accuracy for an ideal unordered-label observation, not a general upper or lower bound on trace inference: actual traces may obscure labels or disclose roles. Saved-prediction analysis now gives exact RF task recognition of 49/99 and 44/99 on the different-label fixtures, and 21/21 and 20/21 on matching-label fixtures. Section~\ref{supp:shadow-diagnostics} separates this conditional task-label calculation from direct first-role and exact-chain endpoints.

\subsection{Achieved timing and the independent clock}
Table~\ref{tab:live-draws} distinguishes sampled holds from achieved record durations. Each row includes all 455 cover calls in that arm/model, including training and validation. Sequential Qwen cover averages a .5198\,s sampled gap and .5660\,s achieved anchor gap; a .0405\,s hold becomes a .1004\,s observed duration. Concurrent duration increases to .1712\,s. Llama's .0768\,s mean hold becomes .1097\,s sequential and .1234\,s concurrent. Thus the live policy does not instantiate the exact timing-distribution transform used by the receipt-level counterfactual. Its measured inference, overhead and altered chronological context must be evaluated directly.

All 143,920 no-cover and 629,530 cover records match independently reassembled TCP/TLS records, across 21,000 connections and 19,244 call windows. All-record median absolute completion-time errors are 1.07 and 1.33\,ms, with 99th percentiles 16.38 and 21.68\,ms; the maximum is 1.675\,s, including setup and teardown. Table~\ref{tab:live-clock} restricts the comparison to call endpoints: the maximum absolute error is 148.39\,ms, and no episode's call order changes under the packet clock.

The separately labelled packet-clock diagnostic refits only the label-free mixture/removal calculation. Its task scores exactly reproduce the relay-clock scores in all eight cover cells, including 112 and 108 for immediate cover. On those held-out immediate calls it flags 181/182 and 182/182 decoys, with zero false flags among 490 and 634 honest calls. This is a sensitivity check for removal on supplied call windows; P/E remain relay-clock primary analyses, and the study makes no physical-network confidentiality claim.

\begin{table*}[t]\centering\small
\caption{Achieved history-based cover timing, all splits (455 cover calls per row). Means in seconds. An anchor gap is measured from the actual prior completed anchor call to cover record start. Sampled durations are minimum server holds, while achieved durations include record-level transport and scheduling.}
\label{tab:live-draws}
\begin{tabularx}{\textwidth}{@{}lYrrrr@{}}\toprule
Model & Arm & Drawn gap & Achieved anchor gap & Drawn hold & Achieved duration\\\midrule
Qwen & History seq. & 0.5198 & 0.5660 & 0.0405 & 0.1004\\
Qwen & History concurrent & 0.5198 & 0.5772 & 0.0405 & 0.1712\\
Llama & History seq. & 0.5865 & 0.6203 & 0.0768 & 0.1097\\
Llama & History concurrent & 0.5865 & 0.6204 & 0.0768 & 0.1234\\
\bottomrule\end{tabularx}
\end{table*}

\begin{table*}[t]\centering\small
\caption{Call-window clock diagnostics across all 300 fixtures per cell. Relay versus independent packet completion errors, in milliseconds. Values are absolute-error median / 99th percentile / maximum. Every call retains its chronological order under the packet clock. Setup/teardown records are excluded from this table.}
\label{tab:live-clock}
\begin{tabularx}{\textwidth}{@{}lYrrr@{}}\toprule
Model & Arm & Call windows & Start: median / p99 / max & End: median / p99 / max\\\midrule
Qwen & None & 1218 & 0.58 / 7.67 / 18.01 & 0.42 / 5.89 / 9.21 \\
Qwen & Immediate & 1673 & 1.70 / 8.57 / 32.76 & 1.23 / 4.20 / 7.98 \\
Qwen & History seq. & 1671 & 1.46 / 8.03 / 10.62 & 1.07 / 4.36 / 7.02 \\
Qwen & History concurrent & 1672 & 1.74 / 26.06 / 42.48 & 1.36 / 11.96 / 89.31 \\
Qwen & Shadow & 2411 & 2.17 / 57.62 / 148.39 & 1.72 / 25.46 / 116.16 \\
Llama & None & 1536 & 0.84 / 7.55 / 68.03 & 0.62 / 5.27 / 8.95 \\
Llama & Immediate & 1998 & 0.58 / 4.49 / 17.04 & 0.41 / 2.31 / 5.63 \\
Llama & History seq. & 1996 & 0.50 / 4.19 / 9.63 & 0.39 / 2.40 / 4.50 \\
Llama & History concurrent & 2002 & 0.58 / 5.54 / 71.35 & 0.43 / 6.65 / 29.04 \\
Llama & Shadow & 3067 & 0.72 / 12.10 / 99.22 & 0.64 / 9.71 / 89.08 \\
\bottomrule\end{tabularx}
\end{table*}


\subsection{Resource costs, paths and business outcomes}
Table~\ref{tab:live-costs} reports both the primary real-chain latency and a descriptive all-work duration through all planner chains and tool-response receipts. This prevents a concurrent dispatcher from hiding queued work behind real-planner completion. Mean all-work duration is 3.369 and 2.835\,s for concurrent history dispatch, and 5.543 and 4.643\,s for shadow execution. The corresponding no-cover means are 2.912 and 3.332\,s. These are warm-session costs from this machine and execution order, not deployment throughput or a causal speedup estimate.

Whole-connection byte factors are 1.13--1.15 for controller policies and 1.31--1.35 for shadows; restricting accounting to supplied call windows gives 1.26--1.33 and 1.68--1.71. Shadow model-token factors are 2.01 and 1.99. Goal/state outcomes for shadow are 112/113 and 82/67 of 120, against no-cover 113/114 and 82/70. Honest paths exactly match no cover on 115/120 and 103/120 shadow fixtures, so even hidden isolated cover does not guarantee identical planning. No business-equivalence hypothesis is substituted for the prespecified accuracy equivalence test.

The formal dataset contains 3,000 primary episodes plus 600 shadow chains; every issued request has a terminal record and no unknown outcome was replaced. Active collector wall time is 38.8\,min for no cover and 193.1\,min for the four cover arms, including their service initialization and model loads but excluding the intervening pool/audit gate and later analyses. The artifact supplies implementations, frozen plans, timing pools and aggregate results. It supports arithmetic inspection but excludes individual record metadata and predictions, message bodies, captures, keys and weights. Feature reconstruction, refitting and independent per-fixture state replay require local inputs; their checks have separate receipts.

\begin{table*}[t]\centering\small
\caption{Live test-set costs. Real and All are mean seconds from admission to real-planner finish and to all planner/tool-response completion, respectively; both exclude session setup/teardown. Record span is the mean first-to-last tool-record observation. Unique/Call bytes are summed-byte factors versus no cover for complete connections and supplied call windows. Calls count all tool/model calls per fixture; token factors include the shadow planner. These are observed costs under persistent caches and phase order, not isolated causal speed measurements.}
\label{tab:live-costs}
\setlength{\tabcolsep}{3.2pt}
\begin{tabularx}{\textwidth}{@{}lYrrrrrrrr@{}}\toprule
Model & Arm & Real s & All s & Span s & Unique$\times$ & Call$\times$ & Tool calls & Model calls & Tokens$\times$\\\midrule
Qwen & None & 2.912 & 2.912 & 2.272 & 1.00 & 1.00 & 4.08 & 4.08 & 1.00 \\
Qwen & Immediate & 3.575 & 3.575 & 2.899 & 1.15 & 1.33 & 5.60 & 4.08 & 1.00 \\
Qwen & History seq. & 4.280 & 4.280 & 3.611 & 1.15 & 1.32 & 5.59 & 4.08 & 1.00 \\
Qwen & History concurrent & 3.328 & 3.369 & 2.896 & 1.15 & 1.33 & 5.60 & 4.08 & 1.00 \\
Qwen & Shadow & 4.804 & 5.543 & 4.681 & 1.31 & 1.68 & 8.14 & 8.14 & 2.01 \\
\midrule
Llama & None & 3.332 & 3.332 & 2.759 & 1.00 & 1.00 & 5.24 & 5.24 & 1.00 \\
Llama & Immediate & 2.916 & 2.916 & 2.420 & 1.13 & 1.27 & 6.80 & 5.28 & 1.01 \\
Llama & History seq. & 3.822 & 3.822 & 3.339 & 1.13 & 1.26 & 6.80 & 5.28 & 1.01 \\
Llama & History concurrent & 2.791 & 2.835 & 2.448 & 1.13 & 1.27 & 6.83 & 5.31 & 1.02 \\
Llama & Shadow & 3.939 & 4.643 & 4.087 & 1.35 & 1.71 & 10.46 & 10.45 & 1.99 \\
\bottomrule\end{tabularx}
\end{table*}


\subsection{Exploratory time and explicit-position ablations}\label{supp:feature-ablation}
These analyses reuse the completed collection after its original test results were known. The code, transformations and testing families were fixed before computing the added scores; this is an exploratory reanalysis, not a new cohort or a preregistered confirmation. The unchanged Full representation reproduces all sixty original candidate count, individual-prediction, validation-score and selection records for the three controller arms and two planners. No new model request or workload replacement enters the analysis.

No-time replaces timestamps, durations, gaps and wall-time features by fixed values, retaining destinations, bytes, call order and the full candidate family. The nominal gap-detector procedure is retained because it also has destination and first-call features. The pure-gap GMM receives a constant gap feature under this transformation; it is retained for suite identity and does not describe a mixture of real gaps. No-position instead aggregates the task classifier's ordered byte slots by destination, removes explicit call-index, previous-destination and is-last features, and omits the two position-only rules. Timing relations and their sequence remain. It is therefore an explicit-position representation/suite ablation, not an order-free observation.

Every variant refits the original LR/RF and removal/decoding procedures with the original splits and seeds. Validation independently selects size and adaptive attacks in each variant. Family A has eight contrasts: Full adaptive minus No-time adaptive, and No-time adaptive minus Full selected size, in each of two history arms and two planners. Family B has four Full-adaptive-minus-No-position-adaptive contrasts. Each family uses two-sided exact McNemar tests followed by Holm correction. Table~\ref{tab:supp-ablation-tests} gives all discordances and tests. The aggregate artifact retains every candidate's validation and test score, including immediate-arm calibration and all selection names.

Three of the four Full-minus-No-time history contrasts reject after correction; Qwen sequential does not ($p=.0931$). None of the four No-time-minus-Full-size contrasts rejects ($p\ge.544$), and none of the explicit-position contrasts rejects ($p\ge.385$). These are results for the fitted representations and reused fixtures. Non-rejection supplies neither feature irrelevance nor equivalence, and the tests omit refitting and workload-selection uncertainty.
\begin{table*}[t]\centering\footnotesize
\caption{All exploratory ablation tests. Discordant counts are condition-only ($+$) and reference-only ($-$), with 120 paired fixtures. Timing family A contains eight contrasts; position family B contains four. Full size denotes its original validation-selected size attack.}\label{tab:supp-ablation-tests}
\setlength{\tabcolsep}{3pt}
\begin{tabularx}{\textwidth}{@{}lllYrrrrr@{}}\toprule
Family & Planner & Policy & Contrast & $+$ & $-$ & Difference (pp) & Raw $p$ & Holm $p$ \\\midrule
A & Qwen & History, seq. & Full $-$ No-time & 32 & 15 & +14.2 & 0.0186 & 0.0931 \\
A & Qwen & History, seq. & No-time $-$ Full size & 0 & 0 & +0.0 & 1 & 1 \\
A & Qwen & History, conc. & Full $-$ No-time & 42 & 9 & +27.5 & $3.4\!\times\!10^{-6}$ & $2.4\!\times\!10^{-5}$ \\
A & Qwen & History, conc. & No-time $-$ Full size & 0 & 0 & +0.0 & 1 & 1 \\
A & Llama & History, seq. & Full $-$ No-time & 27 & 3 & +20.0 & $8.4\!\times\!10^{-6}$ & $5.1\!\times\!10^{-5}$ \\
A & Llama & History, seq. & No-time $-$ Full size & 13 & 15 & -1.7 & 0.8506 & 1 \\
A & Llama & History, conc. & Full $-$ No-time & 33 & 4 & +24.2 & $1.1\!\times\!10^{-6}$ & $8.7\!\times\!10^{-6}$ \\
A & Llama & History, conc. & No-time $-$ Full size & 19 & 10 & +7.5 & 0.136 & 0.5442 \\
B & Qwen & History, seq. & Full $-$ No-position & 9 & 17 & -6.7 & 0.1686 & 0.5059 \\
B & Qwen & History, conc. & Full $-$ No-position & 13 & 5 & +6.7 & 0.0963 & 0.385 \\
B & Llama & History, seq. & Full $-$ No-position & 12 & 11 & +0.8 & 1 & 1 \\
B & Llama & History, conc. & Full $-$ No-position & 7 & 11 & -3.3 & 0.4807 & 0.9614 \\
\bottomrule\end{tabularx}\end{table*}


\subsection{Actual terminal coverage and recognition strata}\label{supp:terminal}
The diagnostic takes the last issued honest business call in each test trace, regardless of whether its action is correct or its requested goal succeeds. A controller-covered terminal has any cover call starting later; a shadow-covered terminal has a later shadow business call. Later means start order rather than completion. The diagnostic retains a separate category for traces without an honest business call; no such case occurs in the eight displayed cover cells. Table~\ref{tab:supp-terminal} reports actual coverage and the original validation-selected adaptive predictions.

The planned omission of the final reference slot in 79/120 fixtures is not the observed terminal count. Qwen's sequential and concurrent uncovered counts are 83 and 73; Llama's are 107 and 80. These counts describe actual trace order, including planner behavior and dispatch overlap. The covered history strata still have adaptive recognition of 29/37, 39/47, 11/13 and 37/40. Stratification is observational and can depend on action and planner behavior. The artifact gives every attack's strata and explicitly descriptive Fisher tests; no causal placement or general superiority claim is inferred from them.
\begin{table*}[t]\centering\footnotesize
\caption{Exploratory stratification by the last issued honest business call. After means later start time; for shadow, coverage requires a later shadow business call. Requested-goal correctness is not part of this definition. Scores are the saved validation-selected adaptive attack; strata are observational, not randomized position interventions.}\label{tab:supp-terminal}
\setlength{\tabcolsep}{3pt}
\begin{tabularx}{\textwidth}{@{}llrrrYY@{}}\toprule
Planner & Policy & No business call & Uncovered & Covered & Uncovered correct/$n$ & Covered correct/$n$ \\\midrule
Qwen & Immediate & 0 & 83 & 37 & 75/83 & 37/37 \\
Qwen & History, seq. & 0 & 83 & 37 & 64/83 & 29/37 \\
Qwen & History, conc. & 0 & 73 & 47 & 60/73 & 39/47 \\
Qwen & Shadow & 0 & 46 & 74 & 41/46 & 13/74 \\
Llama & Immediate & 0 & 107 & 13 & 96/107 & 12/13 \\
Llama & History, seq. & 0 & 107 & 13 & 93/107 & 11/13 \\
Llama & History, conc. & 0 & 80 & 40 & 69/80 & 37/40 \\
Llama & Shadow & 0 & 57 & 63 & 33/57 & 31/63 \\
\bottomrule\end{tabularx}\end{table*}


\subsection{Exact shadow role and task-label diagnostics}\label{supp:shadow-diagnostics}
The frozen role selector predicts that the first reconstructed chain is real when its saved probability is at least .5. The baseline always predicts the start-order role that is most frequent in TRAIN. It chooses shadow-first for both models and therefore scores 53/120 on the test split, whose realized order is 67 real-first and 53 shadow-first. The selector scores 78/120 and 68/120. Exact paired tests, corrected across the two models, give $p=.00124$ and $.0201$ (Table~\ref{tab:supp-shadow-role}). This compares with that training-defined order baseline; it is not a comparison with every uninformed baseline or a proof of general role recovery. Predicting the first-call role also differs from correctly partitioning and selecting an entire chain.

On different-label fixtures, Qwen's validation-selected last-destination rule predicts the real label 34 times, the shadow label 62 times and neither three times. Llama's selected RF gives 44, 40 and 15. A conditional real-versus-shadow exchangeability test gives Holm $p=.0111$ and $.7436$, with the two selected-model tests forming a separate family. Qwen's asymmetry points toward the shadow label, not improved real-task recognition. Label difficulty, priors, cache/fixture histories and start order may contribute to this comparison, so it is not by itself a causal role-leakage test. The complete task-label table and start-order diagnostics are released as aggregates; individual predictions remain local.
\begin{table*}[t]\centering\footnotesize
\caption{Exploratory shadow diagnostics. First-role prediction uses the frozen $p\ge.5$ selector and is compared with the TRAIN-majority start-order baseline by exact paired testing, Holm across two models. Counts /120 except RF task recognition on the 99 different-label pairs. First-role prediction, exact partition and exact real-chain recovery are distinct endpoints.}\label{tab:supp-shadow-role}
\setlength{\tabcolsep}{3pt}
\begin{tabularx}{\textwidth}{@{}lrrrrrY@{}}\toprule
Planner & First role & Order baseline & Holm $p$ & Exact partition & Exact real chain & RF different-label correct/$n$ \\\midrule
Qwen & 78 & 53 & 0.0012 & 48 & 29 & 49/99 \\
Llama & 68 & 53 & 0.0201 & 59 & 32 & 44/99 \\
\bottomrule\end{tabularx}\end{table*}

